\documentclass[12pt,a4paper]{article}
\usepackage[margin=2.5cm]{geometry}
\usepackage[numbers,sort&compress]{natbib}
\usepackage{amsmath,amssymb,amsfonts,amsthm}
\usepackage{graphicx}
\usepackage{hyperref}
\usepackage{microtype}
\usepackage{doi}
\usepackage{booktabs}
\usepackage{enumitem}
\theoremstyle{remark}
\newtheorem*{remark}{Remark}
\theoremstyle{plain}
\newtheorem{theorem}{Theorem}
\newtheorem{lemma}{Lemma}
\newtheorem{corollary}[theorem]{Corollary}
\newtheorem{proposition}{Proposition}
\newtheorem{hypothesis}{Hypothesis}
\theoremstyle{definition}
\newtheorem{definition}{Definition}

\begin{document}

  \title{Yang--Mills Dynamics from Projective Spectral Entropy:\\
  Vertical Heat-Kernel Variation on the Admissible Fibre}

  \author{J.~Beau\\
  \small Independent Researcher, France\\
  \small \texttt{jerome.beau@cosmochrony.org}}

  \date{2026}

  \maketitle

  \begin{abstract}
    The companion paper~\cite{Beau2026h} showed that the horizontal
    variation of the projective spectral entropy functional
    $S_\Pi[g] = \tfrac{1}{2}\log\det' A_g$ with respect to the base
    metric produces the Einstein tensor as the infrared-dominant response,
    via the Seeley--DeWitt coefficient~$a_2$.
    The present paper carries out the complementary vertical variation:
    extending $A_g$ to a Laplace-type operator
    $A_{g,\mathcal{A}} = -(\nabla^{\mathcal{A}})^2 + E$
    on the associated vector bundle of the admissible principal fibre,
    we compute the Seeley--DeWitt coefficient~$a_4$ for the total
    operator and extract the Yang--Mills term
    $c_{\mathrm{YM}}\int \mathrm{Tr}(F_{\mu\nu}F^{\mu\nu})\sqrt{g}\,
    \mathrm{d}^4 x$.
    Variation of $S_\Pi[g, \mathcal{A}]$ with respect to the admissible
    connection $\mathcal{A}$ then yields the Yang--Mills equations
    $D_\mu F^{a\mu\nu} = 0$ in the current-free sector.
    The induced gauge coupling $g_{\mathrm{YM}}^{-2}
    \sim (12 \cdot 16\pi^2)^{-1}\log(\Lambda/\mu)$ is logarithmic
    in the UV cutoff, while Newton's constant
    $G_N^{-1} \sim \ell_{\mathrm{sp}}^{-2}$ is quadratic.
    This spectral hierarchy is a structural consequence of the
    Seeley--DeWitt expansion and does not require a separate
    input.
    The gauge group $G_\Pi$ is taken as input from
    the spectral admissibility
    programme~\cite{Beau2026q6a,Beau2026a35};
    the $\mathrm{SU}(3)$ sector is unconditional:
    $[H\text{-color}]_{\mathrm{eff}}$ (equality of capacity exponents
    $\delta_c = \delta_{\omega c} = \delta_{\omega^2c}$) and
    $[H\text{-color}]_{\mathrm{pointwise}}$ (exact equality at finite~$q$)
    are both established analytically in~\cite{Beau2026a35,Beau2026a36}.
    The central message is:
    $a_2 \to \text{Einstein}$, $a_4 \to \text{Yang--Mills}$,
    from the single functional~$S_\Pi[g, \mathcal{A}]$.
    This spectral stratification suggests that the natural
    organisational space for unification is spectral and
    variational rather than purely group-theoretic: gravity
    and gauge dynamics are separated by geometric direction
    and heat-kernel order, not by a missing symmetry group.
  \end{abstract}

  \noindent\textbf{Keywords:}
  Yang--Mills equations, heat-kernel expansion,
  Seeley--DeWitt coefficients, spectral entropy,
  induced gauge theory, gauge--gravity hierarchy,
  admissible fibre, projective non-injectivity.

  \newpage
  \tableofcontents

  \newpage

  \section{Introduction}
    \label{sec:introduction}

    The Cosmochrony programme derives physical structure from a static
    relational substrate $\chi$ via a generally non-injective
    projection~$\Pi$~\cite{Beau2026m,Cosmochrony}.
    Two lines of development have reached sufficient maturity to be
    combined.
    The first, pursued in the Gravity paper~\cite{Beau2026h} and
    completed in Gravity~3.0, shows that the horizontal variation
    of the projective spectral entropy functional
    \begin{equation}\label{eq:spectral-entropy}
      S_\Pi[g]
      \;=\;
      \tfrac{1}{2}\log\det'\! A_g
    \end{equation}
    with respect to the base metric $g_{\mu\nu}$ reproduces the
    Einstein--Hilbert term as the infrared-dominant
    contribution~\cite{Beau2026h}.
    The second, pursued in the Q-series from Q6a onward, derives the
    admissible gauge group $G_\Pi$, gauge charges, and the Yang--Mills
    principle from the fibre structure of~$\Pi$~\cite{Beau2026q6a}.

    These two results share a common object, $A_g$, but exploit it
    in orthogonal directions.
    The present paper unifies them by extending $A_g$ to a
    Laplace-type operator $A_{g,\mathcal{A}}$ on the total associated
    bundle of $\Pi$ and computing both the gravitational ($a_2$) and
    the gauge ($a_4$) sectors of the spectral entropy from a single
    functional.

    \subsection{Main result}
      \label{subsec:main-result}

      The central result is the following.

      \begin{theorem}[Yang--Mills from $a_4$]
        \label{thm:main}
        Let $G_\Pi$ be the admissible gauge group derived
        in~\cite{Beau2026q6a,Beau2026a35}, and let
        $A_{g,\mathcal{A}}$ be the Laplace-type operator on the
        associated vector bundle constructed in
        Section~\ref{sec:operator}.
        Then:
        \begin{enumerate}[label=\textup{(\roman*)}]
        \item The Seeley--DeWitt coefficient $a_4(A_{g,\mathcal{A}})$
        contains the Yang--Mills density:
        \[
          a_4(A_{g,\mathcal{A}})
          \;\supset\;
          \frac{1}{16\pi^2}
          \int
          \frac{1}{12}\,
          \mathrm{Tr}(\Omega_{\mu\nu}\Omega^{\mu\nu})
          \sqrt{g}\,\mathrm{d}^4 x\,,
          \quad
          \Omega_{\mu\nu} = F_{\mu\nu}\,.
        \]
        \item The projective spectral entropy contains the
        renormalized Yang--Mills term:
        \[
          S_\Pi[g, \mathcal{A}]
          \;\supset\;
          \frac{1}{12 \cdot 16\pi^2}\,
          \log\!\left(\frac{\Lambda}{\mu}\right)
          \int
          \mathrm{Tr}(F_{\mu\nu}F^{\mu\nu})
          \sqrt{g}\,\mathrm{d}^4 x\,.
        \]
        \item The vertical variation of $S_\Pi$ yields the
        Yang--Mills equations in the current-free sector:
        \[
          \frac{\delta S_\Pi}{\delta \mathcal{A}^a_\nu}
          \;=\; 0
          \quad\Longrightarrow\quad
          D_\mu F^{a\mu\nu} = 0\,.
        \]
        \end{enumerate}
      \end{theorem}

      These three statements follow from standard heat-kernel
      technology~\cite{Vassilevich2003,ParkerToms2009} applied to the
      operator $A_{g,\mathcal{A}}$, once the base--fibre structure and
      the admissibility of $\mathcal{A}$ are established.

    \subsection{Position in the programme}
      \label{subsec:position}

      The logical chain of the present paper is:
      \[
        \Pi
        \;\Rightarrow\;
        P_{G_\Pi}(M, G_\Pi)
        \;\Rightarrow\;
        \nabla^{\mathcal{A}}
        \;\Rightarrow\;
        A_{g,\mathcal{A}}
        = -(\nabla^{\mathcal{A}})^2 + E
        \;\Rightarrow\;
        a_4 \supset \tfrac{1}{12}\,
        \mathrm{Tr}(F^2)
        \;\Rightarrow\;
        D_\mu F^{\mu\nu} = 0\,.
      \]
      The companion paper~\cite{Beau2026h} established the horizontal
      half of the same chain:
      \[
        A_g
        \;\Rightarrow\;
        a_2 \supset R
        \;\Rightarrow\;
        G_{\mu\nu} = 0\,.
      \]
      Together, they show that the functional $S_\Pi[g,\mathcal{A}]$
      produces Einstein gravity and Yang--Mills gauge dynamics as the
      $a_2$ and $a_4$ sectors of a single spectral entropy expansion.

      The gauge-gravity synthesis — including the joint equations of
      motion, the ratio $G_N g_{\mathrm{YM}}^2$, and the
      Eddington--Born--Infeld completion of the gauge sector —
      is deferred to the companion paper Q13.

    \subsection{Organisation}
      \label{subsec:organisation}

      Section~\ref{sec:decomposition} establishes the base--fibre
      decomposition of the projective operator and proves the
      horizontal--vertical decoupling lemma.
      Section~\ref{sec:operator} constructs $A_{g,\mathcal{A}}$ on
      the associated bundle.
      Section~\ref{sec:a4} computes the $a_4$ coefficient and
      identifies the Yang--Mills term.
      Section~\ref{sec:variation} carries out the vertical variation
      and derives the Yang--Mills equations.
      Section~\ref{sec:hierarchy} establishes the induced coupling
      and the gauge--gravity spectral hierarchy.
      Section~\ref{sec:status} collects the epistemic status of each
      result and the analytical status of [H-color].
      Section~\ref{sec:conclusion} concludes.

      \newpage

  \section{Base--fibre decomposition of the projective operator}
    \label{sec:decomposition}

    \subsection{The admissible principal bundle}
      \label{subsec:bundle}

      The non-injective projection $\Pi : \Omega \to O$ of
      Foundation~M~\cite{Beau2026m} carries a natural
      base--fibre structure.
      The effective base space $M$ is the Riemannian manifold derived
      in Q5b~\cite{Beau2026q6b} and Q11~\cite{Beau2026q11}:
      the BFS shell stratification of $\mathrm{Heis}_3(\mathbb{Z}/q\mathbb{Z})$
      converges in the Gromov--Hausdorff sense to a Carnot--Carathéodory
      structure with homogeneous dimension $D_{\mathrm{hom}} = 4$
      whose effective metric closure is Lorentzian,
      as established in~\cite{Beau2026q11}.
      The effective fibre is the admissible gauge group
      $G_\Pi = \mathrm{SU}(3) \times \mathrm{SU}(2) \times \mathrm{U}(1)$
      derived in~\cite{Beau2026q6a,Beau2026a35} as the composite of
      three admissibility fixed points of orbit sizes~1, 2,~and~3.

      \begin{definition}[Admissible principal bundle]
        \label{def:bundle}
        The admissible principal bundle associated with the projection
        $\Pi$ is the principal $G_\Pi$-bundle
        $\pi : P \to M$, with structure group $G_\Pi$ and
        connection~$\mathcal{A}$, constructed as follows.
        The total space $P$ is the fibre product of the three orbit
        sub-bundles identified in~\cite{Beau2026q6a,Beau2026a35}.
        The connection $\mathcal{A} \in \Omega^1(P, \mathfrak{g}_\Pi)$
        is the unique $\mathrm{ad}$-equivariant horizontal distribution
        compatible with the admissibility constraints of
        Foundation~M~\cite{Beau2026m}.
        The curvature $F_{\mu\nu}
        = \partial_\mu \mathcal{A}_\nu
        - \partial_\nu \mathcal{A}_\mu
        + [\mathcal{A}_\mu, \mathcal{A}_\nu]$
        takes values in $\mathfrak{g}_\Pi$.
      \end{definition}

    \subsection{The associated vector bundle}
      \label{subsec:associated}

      Let $\rho : G_\Pi \to \mathrm{GL}(V)$ be a representation of
      $G_\Pi$ on a finite-dimensional complex vector space~$V$.
      The associated vector bundle is
      $\mathcal{V} = P \times_\rho V \to M$.
      Sections of~$\mathcal{V}$ carry the combined action of the base
      geometry (through the Levi-Civita connection~$\nabla^g$) and the
      fibre gauge field (through $\mathcal{A}$).
      The covariant derivative on $\Gamma(\mathcal{V})$ is
      \begin{equation}\label{eq:covariant-derivative}
        \nabla^\mathcal{A}_\mu
        \;=\;
        \partial_\mu
        + \omega_\mu
        + \mathcal{A}_\mu\,,
      \end{equation}
      where $\omega_\mu$ is the spin connection on~$M$.

    \subsection{Horizontal--vertical decoupling}
      \label{subsec:decoupling}

      The following lemma is essential for the independence of the
      gravitational ($a_2$) and Yang--Mills ($a_4$) sectors.

      \begin{lemma}[Horizontal--vertical decoupling at leading order]
        \label{lem:decoupling}
        At fixed admissible projection fibre, variations of the
        base metric $g$ and of the fibre connection $\mathcal{A}$
        define independent horizontal and vertical variations of the
        total Laplace-type operator $A_{g,\mathcal{A}}$.
        Mixed variations
        $\delta^2 A_{g,\mathcal{A}} / \delta g\, \delta\mathcal{A}$
        contribute only at order $a_6$ and higher in the
        Seeley--DeWitt expansion and do not modify the leading $a_2$
        Einstein sector or the leading $a_4$ Yang--Mills sector.
      \end{lemma}

      \begin{proof}
        The Seeley--DeWitt coefficients $a_{2k}(A_{g,\mathcal{A}})$
        are local integrals of invariants of degree~$2k$ in
        derivatives of $g_{\mu\nu}$ and
        $\mathcal{A}_\mu$~\cite{Vassilevich2003}.
        At minimal coupling, the leading gauge-dependent contribution
        to $a_4$ is the pure Yang--Mills invariant
        $\mathrm{Tr}(\Omega_{\mu\nu}\Omega^{\mu\nu})$.
        Mixed gauge--gravity terms of the form $R\,\mathrm{Tr}(F^2)$
        are either absent at this order or reduce to curvature
        renormalizations that are independent of the vertical
        variation $\delta_\mathcal{A}$ and therefore do not affect
        the Yang--Mills sector.
        The $a_2$ coefficient depends only on $g$ and $E$
        (not on $\mathcal{A}$ at minimal coupling), and the $a_4$
        coefficient thus splits cleanly into a purely gravitational
        part and a purely gauge part.
        Horizontal and vertical variations are therefore decoupled
        at the $a_2$ and $a_4$ levels.
      \end{proof}

      \begin{remark}
        The decoupling holds under minimal coupling,
        i.e.\ when $E$ does not depend on~$\mathcal{A}$.
        Curvature couplings of the form
        $E \supset \kappa\,F_{\mu\nu}g^{\mu\rho}g^{\nu\sigma}$
        would introduce mixed $a_4$ terms;
        the admissibility constraints of Foundation~M~\cite{Beau2026m}
        impose minimal coupling as the natural choice at leading order.
      \end{remark}

  \section{Laplace-type operator with admissible fibre connection}
    \label{sec:operator}

    \subsection{Construction of the operator}
      \label{subsec:operator-construction}

      The spectral entropy functional of the Gravity
      paper~\cite{Beau2026h} is extended to the associated
      bundle $\mathcal{V}$ by replacing $A_g$ with the operator:
      \begin{equation}\label{eq:total-operator}
        A_{g,\mathcal{A}}
        \;=\;
        -(\nabla^\mathcal{A})^2 + E\,,
      \end{equation}
      acting on $L^2(M, \mathcal{V})$.
      Here $(\nabla^\mathcal{A})^2
  = g^{\mu\nu}\nabla^\mathcal{A}_\mu \nabla^\mathcal{A}_\nu$
      is the connection Laplacian, and $E \in \mathrm{End}(\mathcal{V})$
      is the endomorphism encoding the admissibility correction:
      \begin{equation}\label{eq:endomorphism}
        E
        \;=\;
        \frac{R}{6}\,\mathrm{id}_\mathcal{V}
        + E_{\mathrm{adm}}\,,
      \end{equation}
      where $R/6\,\mathrm{id}_\mathcal{V}$ is the
      Lichnerowicz--Weitzenb\"ock correction ensuring minimal coupling
      to the scalar curvature, and $E_{\mathrm{adm}}$ encodes the
      admissibility constraints of the projection fibre.

    \subsection{Self-adjointness and non-negativity}
      \label{subsec:selfadjoint}

      \begin{proposition}
        \label{prop:self-adjoint}
        The operator $A_{g,\mathcal{A}}$ is self-adjoint and
        non-negative on $L^2(M, \mathcal{V})$ for any admissible
        connection $\mathcal{A}$ and any $E \geq 0$.
      \end{proposition}

      \begin{proof}
        The connection Laplacian $-(\nabla^\mathcal{A})^2$ is
        self-adjoint and non-negative by standard functional
        analysis~\cite{ParkerToms2009}.
        The endomorphism~$E$ is self-adjoint by construction
        (the Lichnerowicz term is a scalar multiple of the identity,
        and $E_{\mathrm{adm}}$ is symmetric under the admissibility
        constraints of~\cite{Beau2026m}).
        Non-negativity for $E \geq 0$ follows from the Bochner identity.
      \end{proof}

    \subsection{The extended spectral entropy}
      \label{subsec:extended-entropy}

      The spectral entropy functional is extended to:
      \begin{equation}\label{eq:extended-entropy}
        S_\Pi[g, \mathcal{A}]
        \;=\;
        \tfrac{1}{2}\,\log\det'\! A_{g,\mathcal{A}}\,,
      \end{equation}
      where $\det'$ denotes the zeta-regularized determinant
      excluding the zero modes.
      When $\mathcal{A} = 0$ and $E = R/6\,\mathrm{id}$,
      this reduces to the functional of~\cite{Beau2026h}.

  \section{Heat-kernel coefficient \(a_4\) and the Yang--Mills term}
    \label{sec:a4}

    \subsection{Seeley--DeWitt expansion}
      \label{subsec:sdw-expansion}

      For the operator $A_{g,\mathcal{A}}$ of
      Equation~\eqref{eq:total-operator}, the heat-kernel
      expansion as $t \to 0^+$ takes the standard
      form~\cite{Vassilevich2003,ParkerToms2009}:
      \begin{equation}\label{eq:heat-kernel-expansion}
        \mathrm{Tr}\!\left(e^{-t\,A_{g,\mathcal{A}}}\right)
        \;\sim\;
        (4\pi t)^{-2}
        \sum_{k=0}^{\infty} t^k\, a_{2k}(A_{g,\mathcal{A}})\,,
        \quad
        t \to 0^+\,,
      \end{equation}
      where in four dimensions:
      \begin{align}
        a_0
        &= \int \mathrm{tr}(\mathrm{id})\,\sqrt{g}\,\mathrm{d}^4 x\,,
        \label{eq:a0} \\
        a_2
        &= \int \mathrm{tr}\!\left(
                               \frac{R}{6}\,\mathrm{id} - E
        \right)\sqrt{g}\,\mathrm{d}^4 x\,,
        \label{eq:a2}
      \end{align}
      and $a_4$ is given in
      Section~\ref{subsec:a4-full} below.

    \subsection{The \texorpdfstring{$a_2$}{a2} coefficient
    and the Einstein sector}
      \label{subsec:a2-einstein}

      For $E = R/6\,\mathrm{id}_\mathcal{V} + E_{\mathrm{adm}}$,
      the $a_2$ coefficient reduces (at minimal coupling) to:
      \begin{equation}\label{eq:a2-minimal}
        a_2(A_{g,\mathcal{A}})
        \;=\;
        -\int \mathrm{tr}(E_{\mathrm{adm}})\,\sqrt{g}\,\mathrm{d}^4 x\,,
      \end{equation}
      which is independent of the gauge connection $\mathcal{A}$.
      Under the identification of the effective metric established
      in~\cite{Beau2026h}, the variation $\delta_{g} S_\Pi
  = \delta_g[\tfrac{1}{2}\log\det' A_{g,\mathcal{A}}]$
      reproduces the Einstein tensor as the infrared-dominant
      contribution; this is the result of the Gravity
      paper~\cite{Beau2026h} and is not repeated here.
      Lemma~\ref{lem:decoupling} guarantees that this sector
      is not modified by the gauge connection at orders $a_2$
      and~$a_4$.

    \subsection{The \texorpdfstring{$a_4$}{a4} coefficient}
      \label{subsec:a4-full}

      The full $a_4$ coefficient for the operator
      $D = -(\nabla^\mathcal{A})^2 + E$ acting on sections of a
      vector bundle with fibre curvature
      $\Omega_{\mu\nu} = [\nabla^\mathcal{A}_\mu,
      \nabla^\mathcal{A}_\nu]$
      is~\cite{Vassilevich2003}:
      \begin{align}\label{eq:a4-full}
        a_4(D)
        &=
        \frac{1}{16\pi^2}
        \int \mathrm{tr}\!\Bigl[
          \tfrac{1}{12}\,
          \Omega_{\mu\nu}\Omega^{\mu\nu}
          + \tfrac{1}{2} E^2
          - \tfrac{1}{6} R E
          + \tfrac{1}{180}\bigl(
          5 R^2
          - 2 R_{\mu\nu}R^{\mu\nu}
          + 2 R_{\mu\nu\rho\sigma}R^{\mu\nu\rho\sigma}
          \bigr)\,\mathrm{id}
          + \text{total derivatives}
          \Bigr]
        \sqrt{g}\,\mathrm{d}^4 x\,.
      \end{align}
      The terms involving $R$, $R_{\mu\nu}$,
      $R_{\mu\nu\rho\sigma}$ are purely gravitational and
      do not depend on $\mathcal{A}$.
      The term $\tfrac{1}{2}E^2$ contributes to the cosmological
      sector and to the admissibility renormalization.
      The gauge-dependent term is:
      \begin{equation}\label{eq:a4-gauge}
        a_4(D)\big|_{\mathrm{gauge}}
        \;=\;
        \frac{1}{16\pi^2}
        \int
        \frac{1}{12}\,
        \mathrm{Tr}(\Omega_{\mu\nu}\Omega^{\mu\nu})
        \sqrt{g}\,\mathrm{d}^4 x\,,
      \end{equation}
      where $\mathrm{Tr}$ denotes the trace over the fibre
      representation $\rho$ of $G_\Pi$.

    \subsection{Identification of the curvature two-form}
      \label{subsec:curvature-identification}

      The fibre curvature of the admissible connection
      $\mathcal{A}$ coincides with the gauge field strength:
      \begin{equation}\label{eq:curvature-identification}
        \Omega_{\mu\nu}
        \;=\;
        [\nabla^\mathcal{A}_\mu,\, \nabla^\mathcal{A}_\nu]
        \;=\;
        \partial_\mu \mathcal{A}_\nu
        - \partial_\nu \mathcal{A}_\mu
        + [\mathcal{A}_\mu, \mathcal{A}_\nu]
        \;=\;
        F_{\mu\nu}\,.
      \end{equation}
      This identification is standard for connection Laplacians
      on principal bundles and holds for any representation
      $\rho$ of $G_\Pi$.
      Substituting into~\eqref{eq:a4-gauge}:
      \begin{equation}\label{eq:a4-YM}
        a_4(A_{g,\mathcal{A}})\big|_{\mathrm{gauge}}
        \;=\;
        \frac{1}{16\pi^2}
        \int
        \frac{1}{12}\,
        \mathrm{Tr}(F_{\mu\nu}F^{\mu\nu})
        \sqrt{g}\,\mathrm{d}^4 x\,.
      \end{equation}
      This establishes part~(i) of Theorem~\ref{thm:main}.

  \section{Vertical variation and Yang--Mills equations}
    \label{sec:variation}

    \subsection{The spectral entropy and its renormalized form}
      \label{subsec:renormalized}

      The zeta-regularized spectral entropy~\eqref{eq:extended-entropy}
      is related to the heat kernel by:
      \begin{equation}\label{eq:zeta-entropy}
        S_\Pi[g, \mathcal{A}]
        \;=\;
        -\tfrac{1}{2}\,\zeta'_{A_{g,\mathcal{A}}}(0)\,,
      \end{equation}
      where $\zeta_D(s) = \mathrm{Tr}(D^{-s})$ is the spectral
      zeta function.
      After UV regularization with cutoff~$\Lambda$ and
      renormalization scale $\mu$, the gauge-dependent part of the
      entropy acquires the contribution:
      \begin{equation}\label{eq:entropy-YM}
        S_\Pi[g, \mathcal{A}]
        \;\supset\;
        \frac{1}{12 \cdot 16\pi^2}\,
        \log\!\left(\frac{\Lambda}{\mu}\right)
        \int
        \mathrm{Tr}(F_{\mu\nu}F^{\mu\nu})
        \sqrt{g}\,\mathrm{d}^4 x\,,
      \end{equation}
      which establishes part~(ii) of Theorem~\ref{thm:main}.
      The logarithmic form of the UV divergence is the hallmark of
      a renormalizable gauge sector in four dimensions.

    \subsection{Vertical variation}
      \label{subsec:vertical-variation}

      The variation of $S_\Pi[g, \mathcal{A}]$ with respect to the
      gauge connection $\mathcal{A}^a_\mu$ is obtained from
      \eqref{eq:entropy-YM} via:
      \begin{align}\label{eq:vertical-variation}
        \delta_\mathcal{A}
        \left[\mathrm{Tr}(F_{\mu\nu}F^{\mu\nu})\right]
        &=
        2\,\mathrm{Tr}
        \!\left[F^{\mu\nu}\,\delta F_{\mu\nu}\right]
        \notag\\
        &=
        2\,\mathrm{Tr}
        \!\left[F^{\mu\nu}
            (D_\mu \delta\mathcal{A}_\nu
            - D_\nu \delta\mathcal{A}_\mu)
        \right]
        \notag\\
        &=
        4\,\mathrm{Tr}
        \!\left[F^{\mu\nu} D_\mu \delta\mathcal{A}_\nu\right]\,.
      \end{align}
      Integrating by parts (using $\nabla_\mu\sqrt{g} = 0$):
      \begin{equation}\label{eq:integration-by-parts}
        \int \mathrm{Tr}
        \!\left[F^{\mu\nu} D_\mu \delta\mathcal{A}_\nu\right]
        \sqrt{g}\,\mathrm{d}^4 x
        =
        -\int \mathrm{Tr}
        \!\left[(D_\mu F^{\mu\nu})\delta\mathcal{A}_\nu\right]
        \sqrt{g}\,\mathrm{d}^4 x
        + \text{boundary terms}\,.
      \end{equation}
      Therefore:
      \begin{equation}\label{eq:variation-entropy}
        \frac{\delta S_\Pi}{\delta \mathcal{A}^a_\nu}
        \;=\;
        -\frac{1}{3 \cdot 16\pi^2}\,
        \log\!\left(\frac{\Lambda}{\mu}\right)
        \,(D_\mu F^{a\mu\nu})\,\sqrt{g}\,.
      \end{equation}

    \subsection{Yang--Mills equations}
      \label{subsec:YM-equations}

      In the absence of projected matter currents
      (the ``current-free sector''), the variational
      principle $\delta_\mathcal{A} S_\Pi = 0$ yields:
      \begin{equation}\label{eq:YM-equations}
        \boxed{D_\mu F^{a\mu\nu} = 0\,,}
      \end{equation}
      which are the Yang--Mills equations for the gauge group $G_\Pi$
      in vacuo.
      This establishes part~(iii) of Theorem~\ref{thm:main}.

      \begin{remark}
        When projected matter fields $\psi$ are included,
        the variation of the matter coupling term in $S_\Pi$
        contributes a current $J^{a\nu} = \delta S_{\mathrm{matter}}
        / \delta \mathcal{A}^a_\nu$, and the full equations
        become $D_\mu F^{a\mu\nu} = J^{a\nu}$.
        The derivation of the matter current from the projective
        structure lies beyond the scope of the present paper.
      \end{remark}

      \begin{remark}
        The Yang--Mills equations~\eqref{eq:YM-equations} emerge
        from $S_\Pi$ via the $a_4$ coefficient, which is four
        orders of derivatives higher than the $a_2$ Einstein
        sector.
        This is not a deficiency but the correct spectral
        signature: the gauge sector is renormalizable (logarithmic
        coupling) while the gravitational sector is
        non-renormalizable (quadratic coupling).
        The Seeley--DeWitt expansion encodes this hierarchy
        automatically.
      \end{remark}

  \section{Induced gauge coupling and gauge--gravity hierarchy}
    \label{sec:hierarchy}

    \subsection{Induced gauge coupling}
      \label{subsec:induced-coupling}

      Comparing the Yang--Mills term~\eqref{eq:entropy-YM}
      with the standard Yang--Mills action
      $S_{\mathrm{YM}} = -(4g_{\mathrm{YM}}^2)^{-1}
      \int \mathrm{Tr}(F_{\mu\nu}F^{\mu\nu})
      \sqrt{g}\,\mathrm{d}^4 x$ (in Euclidean signature),
      the induced coupling is:
      \begin{equation}\label{eq:induced-coupling}
        g_{\mathrm{YM}}^{-2}
        \;\sim\;
        \frac{1}{12 \cdot 16\pi^2}\,
        \log\!\left(\frac{\Lambda}{\mu}\right)\,.
      \end{equation}
      This logarithmic induction is characteristic of asymptotically
      free gauge theories in four dimensions and is the standard
      result of the induced gauge theory programme~\cite{Sakharov1968}.

    \subsection{Newton's constant from the same spectral cutoff}
      \label{subsec:newton}

      The induced Newton's constant from the companion
      paper~\cite{Beau2026h} is:
      \begin{equation}\label{eq:newton-induced}
        G_N^{-1}
        \;\sim\;
        \frac{\dim V}{16\pi^2}\,
        \ell_{\mathrm{sp}}^{-2}\,,
      \end{equation}
      where $\ell_{\mathrm{sp}}$ is the spectral capacity length
      (the same cutoff scale that appears in the Born--Infeld
      completion of the gravitational
      sector~\cite{Beau2026h,Beau2026c}).
      Here $\dim V = \sum_i \dim V_i$ is the total dimension of the
      representation bundle.

    \subsection{Gauge--gravity spectral hierarchy}
      \label{subsec:hierarchy}

      The ratio of the two couplings is:
      \begin{equation}\label{eq:hierarchy}
        G_N\, g_{\mathrm{YM}}^2
        \;\sim\;
        \frac{\ell_{\mathrm{sp}}^2}{\dim V}
        \cdot
        \log\!\left(\frac{\Lambda}{\mu}\right)^{-1}\,.
      \end{equation}
      For $\Lambda \gg \mu$ (UV cutoff well above the renormalization
      scale) and $\ell_{\mathrm{sp}} \ll 1$ (spectral length below
      the Planck scale), this ratio satisfies
      $G_N g_{\mathrm{YM}}^2 \ll 1$, reproducing the observed
      hierarchy between gravitational and gauge strengths.

      \begin{remark}
        The hierarchy $G_N g_{\mathrm{YM}}^2 \ll 1$ is a structural
        consequence of the Seeley--DeWitt expansion and requires no
        separate fine-tuning.
        It follows from the fact that gravity is induced at order~$a_2$
        (quadratic UV divergence) while gauge dynamics are induced at
        order~$a_4$ (logarithmic UV divergence).
        The two sectors share the same spectral cutoff $\ell_{\mathrm{sp}}$,
        which is the only scale in the problem beyond the
        renormalization group flow.
      \end{remark}

  \section{Epistemic status and relation to the Standard Model gauge group}
    \label{sec:status}

    \subsection{Epistemic status of each result}
      \label{subsec:epistemic-table}

      Table~\ref{tab:epistemic} collects the epistemic status
      of the main results of this paper.

      \begin{table}[h]
        \caption{Epistemic status of the main results.}
        \label{tab:epistemic}
        \centering
        \begin{tabular}{lll}
          \toprule
          Result & Status & Dependence \\
          \midrule
          $a_4 \supset \frac{1}{12}\mathrm{Tr}(\Omega^2)$
          & Standard~\cite{Vassilevich2003}
          & None \\
          $\Omega_{\mu\nu} = F_{\mu\nu}$ on admissible bundle
          & Structural
          & Definition~\ref{def:bundle} \\
          Horizontal--vertical decoupling
          & Theorem-level (Lemma~\ref{lem:decoupling})
          & Minimal coupling \\
          $\delta_\mathcal{A} S_\Pi \to D_\mu F^{\mu\nu}$
          & Theorem-level (Theorem~\ref{thm:main})
          & None \\
          $g_{\mathrm{YM}}^{-2} \sim \log(\Lambda/\mu)$
          & Scheme-dependent normalisation
          & UV regularization \\
          $G_N^{-1} \sim \ell_{\mathrm{sp}}^{-2}$
          vs
          $g_{\mathrm{YM}}^{-2} \sim \log\Lambda$
          & Structural hierarchy
          & Seeley--DeWitt \\
          $G_\Pi = \mathrm{SU}(3)
          \times \mathrm{SU}(2) \times \mathrm{U}(1)$
          & Input from~\cite{Beau2026q6a,Beau2026a35}
          & Proved unconditionally; $[H\text{-color}]_{\mathrm{pointwise}}$
          via single-frequency BI fingerprint
          in~\cite{Beau2026a35,Beau2026a36} \\
          \bottomrule
        \end{tabular}
      \end{table}

    \subsection{Analytical status of [H-color]}
      \label{subsec:hcolor}

      The present paper takes $G_\Pi$ as an input from the
      spectral admissibility sub-programme~\cite{Beau2026q6a,Beau2026a35}.
      The SU(3) sector of $G_\Pi$ is established unconditionally for the
      colour-adapted Cayley graph $\mathrm{Cay}(G_q, S^{(3)}_q)$
      in~\cite{Beau2026a35}.
      For the standard Cayley graph, $[H\text{-color}]_{\mathrm{eff}}$
      (equality of capacity exponents in the large-$q$ limit) is established
      analytically in~\cite{Beau2026a35,Beau2026a36} via the intra-coset
      Vandermonde symmetry (O31) and the modulation-frequency argument (O32).
      Furthermore, $[H\text{-color}]_{\mathrm{pointwise}}$ (exact equality at
      finite~$q$) is proved in O31 (Proposition~4.23,~\cite{Beau2026a35}): each
      O12 fingerprint vector is a single additive character of frequency
      $B = c_1 b_1 + c_2 b_2 + c_3 b_3$, and the orbit dilation
      $B \mapsto \omega B$ is a bijection of $\mathbb{Z}/q\mathbb{Z}$, so
      $\sigma_c(n) = \sigma_{\omega c}(n) = \sigma_{\omega^2 c}(n)$ holds exactly.
      The Theorem~\ref{thm:main} of the present paper holds for any
      compact gauge group $G_\Pi$ and is independent of these results;
      they bear only on the identification of the specific group
      $G_\Pi = \mathrm{SU}(3) \times \mathrm{SU}(2) \times \mathrm{U}(1)$.

      The complete programme-level status is:
      \begin{itemize}
        \item $G_\Pi = \mathrm{SU}(2) \times \mathrm{U}(1)$:
        unconditional~\cite{Beau2026q6a};
        \item $G_\Pi = \mathrm{SU}(3)$: unconditional for the
        colour-adapted graph; unconditional at the pointwise level
        ($[H\text{-color}]_{\mathrm{pointwise}}$ proved) for the standard
        graph (Proposition~4.23 of O31~\cite{Beau2026a35});
        \item The finite-$q$ residual variance $R_{\mathrm{var}}$
        is characterised as a modulation bias of order $O(q^{-1})$,
        not an open hypothesis~\cite{Beau2026a36}.
      \end{itemize}

    \subsection{The SU(3) Yang--Mills uniqueness problem}
      \label{subsec:su3-uniqueness}

      The open problem identified in~\cite{Beau2026a35}~(§9.2)
      is the analogue of the Q6a Yang--Mills derivation for SU(3):
      showing that the SU(3) gauge dynamics (structure constants,
      covariant derivative, Yang--Mills equations) emerge uniquely
      from the colour triplet co-admissibility structure.
      This is distinct from the result of the present paper,
      which derives the Yang--Mills equations for any compact $G_\Pi$
      from the $a_4$ heat-kernel coefficient.
      The SU(3) Yang--Mills uniqueness problem concerns the
      \emph{internal} structure of the SU(3) sector, not the
      spectral derivation of the Yang--Mills action.

    \subsection{What Q13 will add}
      \label{subsec:q13}

      The present paper establishes the two independent sectors:
      \begin{align*}
        \frac{\delta S_\Pi}{\delta g^{\mu\nu}}
        &= c_{\mathrm{EH}}\,G_{\mu\nu} + \cdots
        & \text{(Gravity~3.0,~\cite{Beau2026h}, } a_2 \text{)}\\
        \frac{\delta S_\Pi}{\delta \mathcal{A}^a_\nu}
        &= c_{\mathrm{YM}}\,D_\mu F^{a\mu\nu} + \cdots
        & \text{(present paper, } a_4 \text{)}
      \end{align*}
      The joint variational problem
      $\delta_{g, \mathcal{A}} S_\Pi = 0$,
      the joint Einstein--Yang--Mills equations of motion,
      the gauge-gravity coupling at order~$a_6$,
      and the ratio $G_N / g_{\mathrm{YM}}^2$ as a
      precise quantitative prediction, are deferred to the
      companion paper Q13 (Gauge--Gravity Spectral Synthesis).

    \subsection{Spectral stratification as an organisational
    principle}
      \label{subsec:spectral-stratification}

      \begin{remark}
        The spectral stratification identified in this paper ---
        gravity at $a_2$, gauge at $a_4$ --- provides an
        organisational principle that differs fundamentally from
        purely group-theoretic unification schemes.
        In the latter, gravity and gauge interactions are sought
        at the same algebraic level, typically as different
        representations of a common group $G_{\mathrm{unif}}$.
        The present framework suggests instead that they are
        separated by a geometric direction (horizontal versus
        vertical variation on the admissible bundle) and by
        spectral order (two versus four derivatives in the
        heat-kernel expansion).
        Their distinct UV behaviours --- quadratic divergence
        for gravity, logarithmic divergence for gauge --- are
        then consequences of this geometric separation rather
        than independent empirical inputs.
        In this reading, the difficulty of incorporating gravity
        into the gauge-theoretic framework of the Standard Model
        may reflect not a missing symmetry but a difference of
        spectral order: gravity and gauge dynamics are
        not at the same level of the Seeley--DeWitt hierarchy
        of the admissible operator.
        This observation does not imply that strictly algebraic
        unification schemes are incorrect; it suggests that the
        natural organisational space for unification in the
        present framework is spectral and variational rather
        than primarily group-theoretic.
      \end{remark}

  \section{Conclusion}
    \label{sec:conclusion}

    The main result of this paper is the derivation of the
    Yang--Mills equations $D_\mu F^{a\mu\nu} = 0$ as the vertical
    $a_4$ response of the projective spectral entropy functional
    $S_\Pi[g, \mathcal{A}]$, in exact parallel with the derivation
    of the Einstein equations as the horizontal $a_2$ response
    established in~\cite{Beau2026h}.

    The conceptually central observation is the following.
    The Seeley--DeWitt degree of a coefficient encodes the physical
    type of the coupling it induces:
    \begin{align*}
      a_2 &\;\longrightarrow\; \text{gravity,}
      &&\text{quadratic UV divergence,}
      &&\text{perturbatively non-renormalizable sector,} \\
      a_4 &\;\longrightarrow\; \text{gauge,}
      &&\text{logarithmic UV divergence,}
      &&\text{renormalizable sector.}
    \end{align*}
    As a consequence, the gauge--gravity hierarchy
    $G_N g_{\mathrm{YM}}^2 \ll 1$ is not an empirical input
    but a structural output of the Seeley--DeWitt expansion:
    it reflects the fact that gravity and gauge dynamics live
    at different levels of the spectral hierarchy of the same
    functional $S_\Pi[g, \mathcal{A}]$.

    Both sectors emerge from the same spectral cutoff
    $\ell_{\mathrm{sp}}$ and the same non-injective projection
    principle.
    The gauge group $G_\Pi$ is inherited from the admissibility
    sub-programme~\cite{Beau2026q6a,Beau2026a35};
    the SU(3) sector is unconditional at the pointwise level,
    with $[H\text{-color}]_{\mathrm{pointwise}}$ established
    in~\cite{Beau2026a35}.
    The spectral stratification identified here --- developed in
    Section~\ref{subsec:spectral-stratification} --- suggests
    more broadly that the organisational space for unification
    may be spectral and variational rather than primarily
    group-theoretic.

    The complete synthesis --- joint equations of motion,
    Born--Infeld gauge-gravity completion,
    and the quantitative hierarchy ratio ---
    is addressed in Q13.

    \newpage
    \bibliographystyle{plain}
    \bibliography{cosmochrony-bibliography,external-refs}

\end{document}
